\section{Тест производительности}

Для оценки эффективности алгоритма на основе Z-блоков (Z-функции) было проведено его сравнение с \textbf{наивным алгоритмом} поиска подстроки (Brute Force), а также со стандартной библиотечной функцией C++ \texttt{std::search}.

Тестирование проводилось на наборе из 1 000 000 слов текста и образце из 500 слов. Время замерялось только для этапа вычислений (без учета времени чтения с диска). Было смоделировано два сценария:
\begin{enumerate}
    \item \textbf{Случайный текст (Average case):} текст состоит из случайно сгенерированных слов.
    \item \textbf{Худший случай (Worst case):} текст представляет собой миллион одинаковых слов (например, \texttt{"a"}), а образец состоит из 499 слов \texttt{"a"} и одного слова \texttt{"b"} в конце.
\end{enumerate}

\begin{alltt}
$ ./bench.o < ./data/random_test.in
Text words: 1000000
Pattern words: 500

Z-function time:   0.0134285 sec 
Naive search time: 0.0022372 sec 
std::search time:  0.00179256 sec 

$ ./bench.o < ./data/worst_test.in
Text words: 1000000
Pattern words: 500

Z-function time:   0.0174882 sec 
Naive search time: 1.292 sec 
std::search time:  1.24121 sec 
\end{alltt}

\subsection*{Анализ результатов}

1. \textbf{Случайные данные.} На обычном случайном тексте наивный поиск и \texttt{std::search} работают невероятно быстро (около 0.002 сек), так как несовпадение обнаруживается на первом же слове образца. Z-функция отрабатывает за 0.013 сек — небольшое отставание вызвано константными затратами на конкатенацию массивов и выделение памяти под Z-массив.

2. \textbf{Худший случай.} На вырожденных данных ярко проявляется разница в асимптотической сложности. Наивный алгоритм раз за разом успешно сравнивает 499 слов, спотыкается на 500-м слове, сдвигается в тексте всего на одну позицию и начинает проверки заново. Его сложность деградирует до $O(N \times M)$, а время вырастает почти до 1.3 секунд.
Алгоритм на основе Z-блоков благодаря запоминанию истории совпадений в окне $[L, R]$ не делает откатов назад. Его асимптотика остается строго линейной $O(N+M)$. Время работы составило всего 0.017 сек, что дало \textbf{ускорение более чем в 75 раз} по сравнению со стандартной функцией библиотеки C++.
\pagebreak